\Needspace{4\baselineskip}
\begin{Example}
\label{ex:differential-operator-boundary-domains}
现在讨论 Section~\ref{subsec:boundary-value-problems-adjoints} 中考察过的微分算子. 设 $\Omega$ 是 $\mathbb{R}^m$ 中具有光滑边界的区域, 并给定 $p\in\mathbb{N}$. 设 $B_j(\bm{x},D)$, $j=1,\ldots,p$, 是阶数低于 $2p$ 且对 $\bm{x}\in\partial\Omega$ 有良好定义的微分算子. 定义其定义域为
\begin{equation}\label{eq:differential-operator-boundary-domain}
\mathcal{D}_B(L):=\{u\in C_b^{2p}(\Omega)\mid B(\bm{x},D)u(\bm{x})=0,\bm{x}\in\partial\Omega,j=1,\ldots,p\}.
\end{equation}
然后用下式定义从 $L^2(\Omega)$ 到 $L^2(\Omega)$ 的算子 $(\mathcal{D}_B(L),L)$:
\begin{equation}\label{eq:general-differential-operator-action}
L(\bm{x},D)u(\bm{x}):=\sum_{|\alpha|\leq 2p}a_\alpha(\bm{x})D^\alpha u(\bm{x}).
\end{equation}

\end{Example}

\Needspace{4\baselineskip}
\begin{Example}
\label{ex:laplacian-dirichlet-neumann-domains}
特别地, 考察 \textbf{Laplacian}\[\Delta u:=\dfrac{\partial^2u}{\partial x_1^2}+\cdots+\dfrac{\partial^2u}{\partial x_m^2}\]同样将它看作从 $L^2(\Omega)$ 到 $L^2(\Omega)$ 的算子. (暂且设 $\Omega$ 为有界区域.) 考察该算子的两类定义域, 第一类对应于 Dirichlet 数据:
\begin{equation}\label{eq:laplacian-classical-dirichlet-domain}
\widetilde{\mathcal{D}}_D(\Delta):=\{u\in C_b^2(\Omega)\mid u(\bm{x})=0\text{ for }\bm{x}\in\partial\Omega\};
\end{equation}
第二类对应于 Neumann 数据:
\begin{equation}\label{eq:laplacian-classical-neumann-domain}
\widetilde{\mathcal{D}}_N(\Delta):=\left\{u\in C_b^2(\Omega)\mid\dfrac{\partial u}{\partial n}(\bm{x})=0\text{ for }\bm{x}\in\partial\Omega\right\}.
\end{equation}
目前关于算子的可逆性还不能作太多论断, 只能观察到算子 $(\widetilde{\mathcal{D}}_N(\Delta),\Delta)$ 具有非平凡的零空间: 常函数. 将在 Chapter~\ref{chap:linear-elliptic-equations} 中研究这些算子的可逆性.
\end{Example}
